\section{Discussion}\label{sec:discussion}

\subsection{Applying PIPs to larger systems}\label{sec:discussion-system-size}

When the mathematical representation of PIPs changes from explicit monomial expansions to colored multigraphs, PIP formula generation is freed from the constraints of system size. Online evaluation cost still grows substantially with system size, however, and becomes the true bottleneck limiting the reach of global PIPs without a distance cutoff.

Previously, around 20 atoms approached the soft limit of global PIP evaluation: calculation remained possible but became very slow, and explicit degree-five formulas could reach approximately 10 GB. At around 40 atoms, formula storage could approach the hard limit imposed by the memory of a typical single server. For tensor-network descent, 20 atoms is still a very small system; with many element types, fixed costs such as descent may even exceed the tensor contractions themselves. At around a hundred atoms, tensor-network descent has only just finished warming up.

PIPs are one of the principal approaches to constructing potential energy surfaces for reactions, dissociation, and long-range interactions in gas-phase chemical dynamics~\cite{Braams2009PIP,Qu2018PIPReview,Bowman2025PIPPerspective}. Within the range of gas-phase reactions and small clusters, tensor-network descent removes system size as a major obstacle to global PIP evaluation. Extension to larger condensed-phase systems will still require strategies commonly used in local potentials, such as distance cutoffs and atom-centered descriptions, to be incorporated into the PIP framework~\cite{Allen2021aPIP}.

\subsection{A size-independent bound and color explosion}\label{sec:discussion-pattern-count}

The CMG representation directly shows that, for a fixed number of element colors \(E\), degree cutoff \(d_{\max}\), and contraction-width cutoff, the number of colored multigraphs has a fixed upper bound independent of molecular system size. Adding atoms increases the sums and contraction-tensor dimensions for each pattern, but does not increase the number of possible patterns. This is the size-independent bound. In the absence of color exhaustion, the polynomial count equals this bound. There is a threshold degree \(d\) below which no color exhaustion occurs. At degree \(d\) and above, the number of exhausted patterns is nondecreasing with total edge multiplicity. We call this degree the exhaustion threshold.

One of the most fundamental difficulties in fitting potential energy surfaces is the curse of dimensionality. Classical techniques, such as interpolation and ordinary polynomial expansions, can describe three- or four-atom surfaces efficiently and accurately, but their computational costs quickly become unmanageable as the atom count increases. Modern fitting methods rely more heavily on machine learning to address this difficulty and have extended successfully to large molecules.
PIPs were originally linear bases for direct polynomial expansions, rather than neural-network descriptors. They can be viewed as the result of imposing permutation invariance on a polynomial expansion of the potential and equating the coefficients of many equivalent monomials. Each equivalence class of monomials with equal coefficients forms a PIP basis function. The size-independent bound shows that the number of degrees of freedom, or coefficients to be determined, in a degree-\(d\) expansion has an upper bound independent of system size. This provides a way to alleviate the curse of dimensionality independently of machine-learning techniques such as neural networks. It may be an important structural reason for the effectiveness of PIPs in potential energy surfaces of small and medium-sized molecules~\cite{Braams2009PIP,Xie2010MSA,Jiang2013PIPNN,Moberg2021PIPUnrestricted,Hao2025PIPFIAlgorithms}.

At fixed degree \(d\), however, the basis count grows rapidly with the number of elements \(E\), roughly as \(E^d\) to \(E^{2d}\). We call this color explosion. In practice, systems with more element types can still achieve good results using a similar number of polynomials at lower total degree, although the reason remains unclear. Color explosion prevents a compact universal PIP potential from covering the entire periodic table, a capability available in many current universal potentials.

\subsection{A new perspective on fundamental invariants}\label{sec:discussion-fundamental-invariants}

Fundamental invariants (FIs) are often used directly as neural-network inputs, replacing the complete linear PIP basis with a smaller generating set~\cite{Shao2016FINN,Chen2020FINN2,Fu2023FIReview}. FI descriptors are more compact and efficient, but mathematically they no longer span the potential by linear combinations alone; polynomial combinations are required. FIs are therefore usually combined with neural networks. Intuitively, FIs generally offer better accuracy at a fixed descriptor count, whereas PIPs offer better accuracy at a fixed degree cutoff; FIs have a substantial inference-speed advantage in almost all cases. Beyond six atoms, the cost of traditional online polynomial-descriptor evaluation often begins to exceed that of the associated neural network, and its share of the total cost grows with system size.

In practical calculations, complete FI generating sets can be constructed only for small systems; somewhat larger systems require a degree cutoff. The truncated set generally no longer generates the entire invariant polynomial ring. It guarantees generation only of all PIPs through degree \(d\), while many higher-degree PIPs cannot be obtained from the retained generators. The graph representation reveals the rapid growth of complete FI sets with system size: for a system of 1,000 water molecules, the number of generators far exceeds the number of atoms in the observable universe. Degree truncation in large systems therefore selects only a small fraction of the lowest-degree generators from a set too vast to grasp intuitively. Successful applications also show that effective FI-based potential fitting does not require the chosen generators to generate the entire invariant polynomial ring.

Since practical fitting does not require a complete generating set, generator selection can also account for expressive power within the truncation and computational cost. Just as a linear space has many possible bases, an invariant polynomial algebra can have different irredundant generating sets~\cite{Thiery2000GraphInvariants}. The colored-multigraph language makes it straightforward to see that disconnected PIPs can be generated by polynomial combinations of connected PIPs of no greater degree. In the low-degree range with sufficient color capacity, connected graph orbit sums form a natural irredundant set of algebra generators. Finite-capacity color exhaustion introduces additional relations that allow some generators to be eliminated~\cite{Thiery2000GraphInvariants}. In other words, below the exhaustion threshold, connected graphs constitute one choice of FIs. For large systems, connected graphs are nearly equivalent to such a choice; for small systems, color exhaustion allows some connected graphs to be removed from the generating set.

Connected graphs in the present method need not serve as independent descriptors: they can also be used to construct disconnected PIP basis functions at low cost. Under \hyperref[sec:sloped-truncation]{sloped truncation}, FIs or connected graphs serve both as descriptors and as computational intermediates, seeking to combine FI inference speed with the linear expressive power of selected PIPs. Relative to an FI basis, this increases the input dimension and hence the number of network parameters. The Levenberg--Marquardt (LM) algorithm is currently the main optimizer used for PIP-NN training; increasing the parameter count places a substantial burden on training, while its impact on inference is more modest. Connected graphs thus serve both as compact descriptors and as shared intermediates, providing an adjustable tradeoff among evaluation cost, the linear PIP function space, and model parameter count.
